% Macros comuns aos quatro formatos de exportação (tex/formatos/formato-*.tex), derivadas
% de diagramacao/essencia-de-deus/comum.tex.  Cada mestre define a
% geometria e o tratamento das notas ANTES de % Macros comuns aos quatro formatos de exportação (tex/formatos/formato-*.tex), derivadas
% de diagramacao/essencia-de-deus/comum.tex.  Cada mestre define a
% geometria e o tratamento das notas ANTES de % Macros comuns aos quatro formatos de exportação (tex/formatos/formato-*.tex), derivadas
% de diagramacao/essencia-de-deus/comum.tex.  Cada mestre define a
% geometria e o tratamento das notas ANTES de % Macros comuns aos quatro formatos de exportação (tex/formatos/formato-*.tex), derivadas
% de diagramacao/essencia-de-deus/comum.tex.  Cada mestre define a
% geometria e o tratamento das notas ANTES de \input{comum}; aqui ficam fonte, títulos,
% marcas de página e cabeçalhos.
\usepackage{fontspec}
\usepackage{polyglossia}
\setmainlanguage{portuguese}
\setotherlanguage{latin}
\setmainfont{EB Garamond}[Numbers={OldStyle,Proportional}]
\usepackage{microtype}
\usepackage{xcolor}
\definecolor{marca}{gray}{0.45}
\setlength{\emergencystretch}{1.5em}
\frenchspacing

% --- página do original ------------------------------------------------------
% \pg{443}: virada de página do original no corpo -> na margem externa.
% \pgn{444}: virada dentro de nota (marginpar é proibido ali) -> discreta no fluxo.
% Em modo vertical (antes de um título ou parágrafo) a marca espera o próximo
% parágrafo, senão cairia no pé da página anterior.
% marginnote, não \marginpar: \marginpar é float e, com notas retidas de uma página
% para outra (\dimen\footins limitado), o LaTeX aborta com «Float(s) lost».
\usepackage{marginnote}
\renewcommand*{\marginfont}{\footnotesize\color{marca}}
\newcommand{\margempg}[1]{\marginnote{[#1]}}
% Antes de \tese/\teseB/\subp (que fazem \par\medskip\noindent) a marca fica pendente
% e é solta por eles (\soltapg) já dentro do parágrafo novo; antes de texto comum,
% abre o parágrafo ali mesmo.  (Um hook para/begin com \marginnote comia o recuo de
% parágrafo.)
\gdef\pgpendente{}
\newcommand{\soltapg}{\pgpendente\gdef\pgpendente{}}
\newcommand{\pg}[1]{%
  \ifvmode\def\pgarg{#1}\expandafter\pgvert
  \else\margempg{#1}\expandafter\ignorespaces\fi}
\newcommand{\pgvert}{\futurelet\pgprox\pgvertdecide}
\newcommand{\pgvertdecide}{%
  \ifx\pgprox\tese \pgadia\else\ifx\pgprox\sub \pgadia\else\ifx\pgprox\teseB \pgadia\else\ifx\pgprox\subp \pgadia
  \else\leavevmode\margempg{\pgarg}\fi\fi\fi\fi}
\newcommand{\pgadia}{%
  \xdef\pgpendente{\unexpanded\expandafter{\pgpendente}\noexpand\margempg{\pgarg}}}
\newcommand{\pgn}[1]{{\scriptsize\color{marca}[#1]}}
% (classe book, twoside: \marginpar vai para a margem externa)
\setlength{\marginparsep}{3mm}

% --- estrutura do autor (mesmo dialeto de tex/preambulo.tex e de tools/tei2latex.py) ---
% tese numerada que corre com o texto / em linha própria / subdivisão a), b)
\newcommand{\tese}[2]{\par\medskip\noindent\soltapg#1.\enspace\textbf{#2}\;---\;}
\newcommand{\teseB}[2]{\par\medskip\noindent\soltapg#1.\enspace\textbf{#2}\par}
\newcommand{\sub}[2]{\par\medskip\noindent\soltapg#1)\enspace\emph{#2}\,---\;}
\newcommand{\subp}[1]{\par\noindent\soltapg\hspace*{\parindent}#1)\enspace}
% asterisco do original: abre um apêndice sem número dentro de uma nota
\newcommand{\asterisco}{\par\smallskip{\centering$\ast$\par}\smallskip\noindent}

% gravura recortada do fac-símile (<figure> do TEI)
\usepackage{graphicx}
\newcommand{\gravura}[3]{\par\medskip\begingroup\centering
  \includegraphics[width=#2\linewidth,height=0.75\textheight,keepaspectratio]{#1}\par
  \if\relax\detokenize{#3}\relax\else\smallskip{\small #3\par}\fi\endgroup\medskip\noindent}
\newenvironment{latim}{\begin{quote}\itshape}{\end{quote}}

% --- cabeçalho: título corrente (\def\titulocorrente{...} no -usepretex) ---------
\providecommand{\titulocorrente}{elementos de philosophia}
\usepackage{fancyhdr}
\pagestyle{fancy}
\fancyhf{}
\fancyhead[LE]{\small\oldstylenums{\thepage}\quad\scshape\titulocorrente}
\fancyhead[RO]{\small\scshape\titulocorrente\quad\oldstylenums{\thepage}}
\renewcommand{\headrulewidth}{0pt}
\fancypagestyle{plain}{\fancyhf{}\fancyfoot[C]{\small\oldstylenums{\thepage}}}
\usepackage{sectsty}\allsectionsfont{\raggedright}
\providecommand{\fragmento}{build/vol2}
; aqui ficam fonte, títulos,
% marcas de página e cabeçalhos.
\usepackage{fontspec}
\usepackage{polyglossia}
\setmainlanguage{portuguese}
\setotherlanguage{latin}
\setmainfont{EB Garamond}[Numbers={OldStyle,Proportional}]
\usepackage{microtype}
\usepackage{xcolor}
\definecolor{marca}{gray}{0.45}
\setlength{\emergencystretch}{1.5em}
\frenchspacing

% --- página do original ------------------------------------------------------
% \pg{443}: virada de página do original no corpo -> na margem externa.
% \pgn{444}: virada dentro de nota (marginpar é proibido ali) -> discreta no fluxo.
% Em modo vertical (antes de um título ou parágrafo) a marca espera o próximo
% parágrafo, senão cairia no pé da página anterior.
% marginnote, não \marginpar: \marginpar é float e, com notas retidas de uma página
% para outra (\dimen\footins limitado), o LaTeX aborta com «Float(s) lost».
\usepackage{marginnote}
\renewcommand*{\marginfont}{\footnotesize\color{marca}}
\newcommand{\margempg}[1]{\marginnote{[#1]}}
% Antes de \tese/\teseB/\subp (que fazem \par\medskip\noindent) a marca fica pendente
% e é solta por eles (\soltapg) já dentro do parágrafo novo; antes de texto comum,
% abre o parágrafo ali mesmo.  (Um hook para/begin com \marginnote comia o recuo de
% parágrafo.)
\gdef\pgpendente{}
\newcommand{\soltapg}{\pgpendente\gdef\pgpendente{}}
\newcommand{\pg}[1]{%
  \ifvmode\def\pgarg{#1}\expandafter\pgvert
  \else\margempg{#1}\expandafter\ignorespaces\fi}
\newcommand{\pgvert}{\futurelet\pgprox\pgvertdecide}
\newcommand{\pgvertdecide}{%
  \ifx\pgprox\tese \pgadia\else\ifx\pgprox\sub \pgadia\else\ifx\pgprox\teseB \pgadia\else\ifx\pgprox\subp \pgadia
  \else\leavevmode\margempg{\pgarg}\fi\fi\fi\fi}
\newcommand{\pgadia}{%
  \xdef\pgpendente{\unexpanded\expandafter{\pgpendente}\noexpand\margempg{\pgarg}}}
\newcommand{\pgn}[1]{{\scriptsize\color{marca}[#1]}}
% (classe book, twoside: \marginpar vai para a margem externa)
\setlength{\marginparsep}{3mm}

% --- estrutura do autor (mesmo dialeto de tex/preambulo.tex e de tools/tei2latex.py) ---
% tese numerada que corre com o texto / em linha própria / subdivisão a), b)
\newcommand{\tese}[2]{\par\medskip\noindent\soltapg#1.\enspace\textbf{#2}\;---\;}
\newcommand{\teseB}[2]{\par\medskip\noindent\soltapg#1.\enspace\textbf{#2}\par}
\newcommand{\sub}[2]{\par\medskip\noindent\soltapg#1)\enspace\emph{#2}\,---\;}
\newcommand{\subp}[1]{\par\noindent\soltapg\hspace*{\parindent}#1)\enspace}
% asterisco do original: abre um apêndice sem número dentro de uma nota
\newcommand{\asterisco}{\par\smallskip{\centering$\ast$\par}\smallskip\noindent}

% gravura recortada do fac-símile (<figure> do TEI)
\usepackage{graphicx}
\newcommand{\gravura}[3]{\par\medskip\begingroup\centering
  \includegraphics[width=#2\linewidth,height=0.75\textheight,keepaspectratio]{#1}\par
  \if\relax\detokenize{#3}\relax\else\smallskip{\small #3\par}\fi\endgroup\medskip\noindent}
\newenvironment{latim}{\begin{quote}\itshape}{\end{quote}}

% --- cabeçalho: título corrente (\def\titulocorrente{...} no -usepretex) ---------
\providecommand{\titulocorrente}{elementos de philosophia}
\usepackage{fancyhdr}
\pagestyle{fancy}
\fancyhf{}
\fancyhead[LE]{\small\oldstylenums{\thepage}\quad\scshape\titulocorrente}
\fancyhead[RO]{\small\scshape\titulocorrente\quad\oldstylenums{\thepage}}
\renewcommand{\headrulewidth}{0pt}
\fancypagestyle{plain}{\fancyhf{}\fancyfoot[C]{\small\oldstylenums{\thepage}}}
\usepackage{sectsty}\allsectionsfont{\raggedright}
\providecommand{\fragmento}{build/vol2}
; aqui ficam fonte, títulos,
% marcas de página e cabeçalhos.
\usepackage{fontspec}
\usepackage{polyglossia}
\setmainlanguage{portuguese}
\setotherlanguage{latin}
\setmainfont{EB Garamond}[Numbers={OldStyle,Proportional}]
\usepackage{microtype}
\usepackage{xcolor}
\definecolor{marca}{gray}{0.45}
\setlength{\emergencystretch}{1.5em}
\frenchspacing

% --- página do original ------------------------------------------------------
% \pg{443}: virada de página do original no corpo -> na margem externa.
% \pgn{444}: virada dentro de nota (marginpar é proibido ali) -> discreta no fluxo.
% Em modo vertical (antes de um título ou parágrafo) a marca espera o próximo
% parágrafo, senão cairia no pé da página anterior.
% marginnote, não \marginpar: \marginpar é float e, com notas retidas de uma página
% para outra (\dimen\footins limitado), o LaTeX aborta com «Float(s) lost».
\usepackage{marginnote}
\renewcommand*{\marginfont}{\footnotesize\color{marca}}
\newcommand{\margempg}[1]{\marginnote{[#1]}}
% Antes de \tese/\teseB/\subp (que fazem \par\medskip\noindent) a marca fica pendente
% e é solta por eles (\soltapg) já dentro do parágrafo novo; antes de texto comum,
% abre o parágrafo ali mesmo.  (Um hook para/begin com \marginnote comia o recuo de
% parágrafo.)
\gdef\pgpendente{}
\newcommand{\soltapg}{\pgpendente\gdef\pgpendente{}}
\newcommand{\pg}[1]{%
  \ifvmode\def\pgarg{#1}\expandafter\pgvert
  \else\margempg{#1}\expandafter\ignorespaces\fi}
\newcommand{\pgvert}{\futurelet\pgprox\pgvertdecide}
\newcommand{\pgvertdecide}{%
  \ifx\pgprox\tese \pgadia\else\ifx\pgprox\sub \pgadia\else\ifx\pgprox\teseB \pgadia\else\ifx\pgprox\subp \pgadia
  \else\leavevmode\margempg{\pgarg}\fi\fi\fi\fi}
\newcommand{\pgadia}{%
  \xdef\pgpendente{\unexpanded\expandafter{\pgpendente}\noexpand\margempg{\pgarg}}}
\newcommand{\pgn}[1]{{\scriptsize\color{marca}[#1]}}
% (classe book, twoside: \marginpar vai para a margem externa)
\setlength{\marginparsep}{3mm}

% --- estrutura do autor (mesmo dialeto de tex/preambulo.tex e de tools/tei2latex.py) ---
% tese numerada que corre com o texto / em linha própria / subdivisão a), b)
\newcommand{\tese}[2]{\par\medskip\noindent\soltapg#1.\enspace\textbf{#2}\;---\;}
\newcommand{\teseB}[2]{\par\medskip\noindent\soltapg#1.\enspace\textbf{#2}\par}
\newcommand{\sub}[2]{\par\medskip\noindent\soltapg#1)\enspace\emph{#2}\,---\;}
\newcommand{\subp}[1]{\par\noindent\soltapg\hspace*{\parindent}#1)\enspace}
% asterisco do original: abre um apêndice sem número dentro de uma nota
\newcommand{\asterisco}{\par\smallskip{\centering$\ast$\par}\smallskip\noindent}

% gravura recortada do fac-símile (<figure> do TEI)
\usepackage{graphicx}
\newcommand{\gravura}[3]{\par\medskip\begingroup\centering
  \includegraphics[width=#2\linewidth,height=0.75\textheight,keepaspectratio]{#1}\par
  \if\relax\detokenize{#3}\relax\else\smallskip{\small #3\par}\fi\endgroup\medskip\noindent}
\newenvironment{latim}{\begin{quote}\itshape}{\end{quote}}

% --- cabeçalho: título corrente (\def\titulocorrente{...} no -usepretex) ---------
\providecommand{\titulocorrente}{elementos de philosophia}
\usepackage{fancyhdr}
\pagestyle{fancy}
\fancyhf{}
\fancyhead[LE]{\small\oldstylenums{\thepage}\quad\scshape\titulocorrente}
\fancyhead[RO]{\small\scshape\titulocorrente\quad\oldstylenums{\thepage}}
\renewcommand{\headrulewidth}{0pt}
\fancypagestyle{plain}{\fancyhf{}\fancyfoot[C]{\small\oldstylenums{\thepage}}}
\usepackage{sectsty}\allsectionsfont{\raggedright}
\providecommand{\fragmento}{build/vol2}
; aqui ficam fonte, títulos,
% marcas de página e cabeçalhos.
\usepackage{fontspec}
\usepackage{polyglossia}
\setmainlanguage{portuguese}
\setotherlanguage{latin}
\setmainfont{EB Garamond}[Numbers={OldStyle,Proportional}]
\usepackage{microtype}
\usepackage{xcolor}
\definecolor{marca}{gray}{0.45}
\setlength{\emergencystretch}{1.5em}
\frenchspacing

% --- página do original ------------------------------------------------------
% \pg{443}: virada de página do original no corpo -> na margem externa.
% \pgn{444}: virada dentro de nota (marginpar é proibido ali) -> discreta no fluxo.
% Em modo vertical (antes de um título ou parágrafo) a marca espera o próximo
% parágrafo, senão cairia no pé da página anterior.
% marginnote, não \marginpar: \marginpar é float e, com notas retidas de uma página
% para outra (\dimen\footins limitado), o LaTeX aborta com «Float(s) lost».
\usepackage{marginnote}
\renewcommand*{\marginfont}{\footnotesize\color{marca}}
\newcommand{\margempg}[1]{\marginnote{[#1]}}
% Antes de \tese/\teseB/\subp (que fazem \par\medskip\noindent) a marca fica pendente
% e é solta por eles (\soltapg) já dentro do parágrafo novo; antes de texto comum,
% abre o parágrafo ali mesmo.  (Um hook para/begin com \marginnote comia o recuo de
% parágrafo.)
\gdef\pgpendente{}
\newcommand{\soltapg}{\pgpendente\gdef\pgpendente{}}
\newcommand{\pg}[1]{%
  \ifvmode\def\pgarg{#1}\expandafter\pgvert
  \else\margempg{#1}\expandafter\ignorespaces\fi}
\newcommand{\pgvert}{\futurelet\pgprox\pgvertdecide}
\newcommand{\pgvertdecide}{%
  \ifx\pgprox\tese \pgadia\else\ifx\pgprox\sub \pgadia\else\ifx\pgprox\teseB \pgadia\else\ifx\pgprox\subp \pgadia
  \else\leavevmode\margempg{\pgarg}\fi\fi\fi\fi}
\newcommand{\pgadia}{%
  \xdef\pgpendente{\unexpanded\expandafter{\pgpendente}\noexpand\margempg{\pgarg}}}
\newcommand{\pgn}[1]{{\scriptsize\color{marca}[#1]}}
% (classe book, twoside: \marginpar vai para a margem externa)
\setlength{\marginparsep}{3mm}

% --- estrutura do autor (mesmo dialeto de tex/preambulo.tex e de tools/tei2latex.py) ---
% tese numerada que corre com o texto / em linha própria / subdivisão a), b)
\newcommand{\tese}[2]{\par\medskip\noindent\soltapg#1.\enspace\textbf{#2}\;---\;}
\newcommand{\teseB}[2]{\par\medskip\noindent\soltapg#1.\enspace\textbf{#2}\par}
\newcommand{\sub}[2]{\par\medskip\noindent\soltapg#1)\enspace\emph{#2}\,---\;}
\newcommand{\subp}[1]{\par\noindent\soltapg\hspace*{\parindent}#1)\enspace}
% asterisco do original: abre um apêndice sem número dentro de uma nota
\newcommand{\asterisco}{\par\smallskip{\centering$\ast$\par}\smallskip\noindent}

% gravura recortada do fac-símile (<figure> do TEI)
\usepackage{graphicx}
\newcommand{\gravura}[3]{\par\medskip\begingroup\centering
  \includegraphics[width=#2\linewidth,height=0.75\textheight,keepaspectratio]{#1}\par
  \if\relax\detokenize{#3}\relax\else\smallskip{\small #3\par}\fi\endgroup\medskip\noindent}
\newenvironment{latim}{\begin{quote}\itshape}{\end{quote}}

% --- cabeçalho: título corrente (\def\titulocorrente{...} no -usepretex) ---------
\providecommand{\titulocorrente}{elementos de philosophia}
\usepackage{fancyhdr}
\pagestyle{fancy}
\fancyhf{}
\fancyhead[LE]{\small\oldstylenums{\thepage}\quad\scshape\titulocorrente}
\fancyhead[RO]{\small\scshape\titulocorrente\quad\oldstylenums{\thepage}}
\renewcommand{\headrulewidth}{0pt}
\fancypagestyle{plain}{\fancyhf{}\fancyfoot[C]{\small\oldstylenums{\thepage}}}
\usepackage{sectsty}\allsectionsfont{\raggedright}
\providecommand{\fragmento}{build/vol2}
